\section{从现实决策到学习问题}
\label{sec:13-Machine-Learning/01-Learning-Problems-and-Evaluation/01}

一家医院攒下了三十万张胸片和对应的诊断结论，数据团队接到的需求只有一句话：训练一个模型预测患者有没有病。听上去边界清楚。可是同一句话至少对应四件不同的事：体检筛查时尽量不漏人，门诊确诊前减少不必要的增强扫描，给放射科医生排一个当天优先阅片的队列，估计某位患者未来一年内发病的概率。四件事共用同一张数据表，但输出的类型不同、错判的代价不同、模型该学的统计量也不同。需求方说的是同一句中文，工程上要写成四个学习问题。

写歪的地方几乎从不出现在优化器里。它出现在从现实处境到数学对象的那一次翻译。

\subsection{数据表不决定学习问题，行动才决定}

先把整条链条摆出来，再逐段检查。现实处境经过特征提取变成 \(x\)，模型给出预测 \(f(x)\)，某条决策规则把预测变成行动 \(a\)，行动落到真实结果 \(y\) 上产生代价。链条上的每个箭头都是一次翻译，每次翻译都可能把问题换成另一个问题。

\begin{center}
\begin{tikzpicture}[>=stealth, every node/.style={font=\small}]

\foreach \x/\lab in {0/{现实处境\\\(x\)}, 3.4/{预测\\\(f(x)\)}, 6.8/{行动\\\(a=d(f(x))\)}, 10.2/{结果与代价\\\(L(a,y)\)}}
  \node[draw, rounded corners, minimum width=2.4cm, minimum height=1.1cm, align=center] at (\x,0) {\lab};

\foreach \x in {0,3.4,6.8}
  \draw[->, thick] (\x+1.3,0) -- (\x+2.1,0);

\node[align=center, font=\scriptsize] at (1.7,0.78) {选预测单位\\与信息时点};
\node[align=center, font=\scriptsize] at (5.1,0.78) {选阈值\\与决策规则};
\node[align=center, font=\scriptsize] at (8.5,0.78) {写下代价\\是否对称};

\draw[->, dashed] (10.2,-0.75) -- (10.2,-1.6) -- (0,-1.6) -- (0,-0.75);
\node[font=\scriptsize] at (5.1,-1.85) {行动改变了下一批能被观察到的数据};

\end{tikzpicture}
\end{center}

\emph{四个方框住在不同的空间，三个箭头各自是一次可能写歪的翻译。}

多数团队的注意力落在中间那个方框上，而失败通常发生在箭头上。凸优化里有一句同源的告诫，见\hyperref[sec:09-Convex-Optimization/01-Optimization-Foundations/01]{``把现实选择翻译成优化问题''}：目标写歪了，解会精确地钻那个空子。学习问题的处境更糟一层——优化器不但会钻空子，还会给你一条漂亮的下降曲线作为掩护。

\subsection{预测单位和信息时点必须先于模型固定下来}

设一次预测所针对的对象为 \(u\)。它可以是一位患者、一张影像、一次交易，也可以是某个用户在某个时刻的状态。特征 \(X\) 只允许包含作出预测那一刻真正拿得到的信息，标签 \(Y\) 描述希望预测的结果。

这两句话看着朴素，却直接改写样本的概率结构。同一位患者的十张影像不是十个独立个体；同一台设备连续采集的记录共享一份噪声；预测住院结局时，出院小结里才写入的诊断编码不能进入输入。把数据库的一行机械地当成一个独立样本，等于在问题定义阶段凭空制造样本量，或者把未来的信息偷渡进现在。

所以数据要写成随机对象，而不只是一个矩阵：

\[
Z=(X,Y)\sim P .
\]

这里的 \(P\) 不只描述数值分布，也隐含了预测单位、采样方式和时间边界。写下 \(n\) 个样本

\[
S=(Z_1,\ldots,Z_n)
\]

之后，常用的独立同分布假设 \(Z_i\overset{\iid}{\sim}P\) 是一项需要逐条核对的建模承诺，不是表格有 \(n\) 行就自动到手的性质。

\subsection{输出、行动和结果是三个对象}

模型可能输出分数 \(s(x)\in\R\)、概率 \(p(x)\in[0,1]\)、类别 \(\hat y(x)\)，或者一段表示向量。现实系统最终执行的却是行动 \(a\in\mathcal A\)。概率模型给出患病概率，医院再根据复查容量和漏诊代价决定要不要召回。上图中间那两个方框之所以要分开画，正是因为二者常被当成一件事。

把它们混在一起会长出几种典型误解。一个概率估计器可以校准得很好，却因为阈值取在 \(1/2\) 而造成糟糕决策；一个只保证排序的分数可以有漂亮的 AUC，却不能当成概率来读；同一个模型放到复查容量和错误成本都不同的两家医院，也不该沿用同一个阈值。数理统计中的\hyperref[ch:092]{``统计决策与可信预测：模型给出概率之后怎样行动''}把这条接口完整展开：后验概率怎样经由损失变成 Bayes 行动，校准为何不等于区分能力，覆盖率保证又为何不等于每个个体都有条件保证。

现实结果由损失函数记账：

\[
L(a,y),
\]

它衡量真实结果为 \(y\) 时采取行动 \(a\) 的代价。先说清行动和代价，``预测得好''才成为一句可以被推翻的话。

\subsection{假设空间写下允许怎样从输入走到输出}

学习算法不会在所有可能函数里毫无限制地搜索。它先选定一个假设空间

\[
\mathcal H=\set{f_\theta:\theta\in\Theta},
\]

再用数据挑参数。线性回归把 \(\mathcal H\) 限制成仿射函数，决策树限制成递归分区，核方法通过再生核 Hilbert 空间的范数控制函数的振荡，神经网络则用层数、宽度、激活和参数共享规定可表达的复合结构。

这种限制就是归纳偏置。数据只约束已经观察到的位置，模型怎样延伸到没观察过的位置，由 \(\mathcal H\) 的结构决定。两个模型都能把训练误差压到零，在训练点之间仍可能给出完全不同的函数——而部署时遇到的点，多半正落在那些缝隙里。

参数化也不等于模型本身。不同参数可以表示同一个函数：隐藏单元置换后网络输出不变，同一函数族常有多种坐标表示。讨论可识别性之前，要先分清想识别的是参数、预测函数，还是某个现实机制。

\subsection{训练分布与部署分布必须分开命名}

历史数据来自 \(P_{\mathrm{train}}\)，真正关心的未来对象来自 \(P_{\mathrm{deploy}}\)。教科书常暂时假设两者相同：

\[
P_{\mathrm{train}}=P_{\mathrm{deploy}} .
\]

这个等式是独立测试集能够模拟部署表现的全部依据，也往往是应用中最脆弱的一环。医院换了 CT 机、平台改了推荐策略、欺诈者摸清了检测规则，都会改变输入分布、标签分布，或者二者之间的关系。

真正该问的是样本是否来自你想下结论的那个总体。一百万条旧设备样本替代不了一千条新设备样本；样本量只能压低对既定分布的抽样误差，修不了目标总体选错造成的偏差。上图那条虚线还提示了更麻烦的一种情形：行动本身参与生成下一批数据，此时 \(P_{\mathrm{deploy}}\) 是被模型改写过的分布。

\subsection{三种范式改变的是数据从哪里来}

监督学习观察输入与目标的配对样本，学一条预测规则；无监督学习没有外部标签，目标来自密度、重构、聚类或表示上的约束；强化学习中的行动会改变后续能观察到的数据，样本分布反过来依赖当前策略。

这三者不能只用``有没有标签''分开。无监督方法照样在优化一个人选的目标，聚类结果并不是数据里唯一存在的真分类；强化学习也可能手握带奖励的样本，却因为数据由行为策略产生而不能当成普通监督数据来评估。范式决定了哪些概率对象可被观察，也决定了评估时需要模拟怎样的数据生成过程。

\subsection{把问题写成可审计的六句话}

在挑算法之前，学习问题应该已经被写成六句能被别人反驳的话：预测单位是谁，同一单位会不会贡献多行相关记录；预测发生在哪个时刻，那一刻哪些变量真正可用；模型输出的是数值、概率、排序、类别还是行动；错误怎样传导到现实决策，两类代价是否对称；训练样本经由怎样的机制被收集，目标部署总体又是谁；最终希望识别参数、预测未来，还是解释因果机制。

最后一句尤其容易含混。预测模型利用相关结构，并不自动识别干预效果。若治疗选择本身受病情严重程度影响，数据里``接受治疗的人结局更差''可能完全来自选择偏差，而不是治疗有害。再高的预测精度也无法把条件相关性改写成因果结论，那需要另外一套干预、混杂与可识别性假设。

六句话写完，问题才算立住。接下来要处理的是它内部的第二重混淆：单样本的代价、总体上的期望代价、训练时实际在跑的表达式、汇报时端出去的那个数字，这四个量在代码里只隔几行，在数学上并不住在一起。下一节从一份风控日志开始拆它们。
